\subsection{Determinante von $2\times2$-Matrizen}
\label{subsec:linear_systems_two_by_two_determinants}

Wir lösen nun das allgemeine $2\times2$-LGS mit dem Gauß-Verfahren, um die
Struktur seiner Lösung genauer zu untersuchen.

\begin{ExampleBox}{Allgemeines $2\times2$-LGS}{linear_systems_determinant_from_gaussian_elimination}
  Wir lösen das folgende $2\times2$-LGS mit dem Gauß-Verfahren
  \[
    \begin{aligned}
      a_{11}x_1+a_{12}x_2&=b_1,\\
      a_{21}x_1+a_{22}x_2&=b_2
    \end{aligned}
    \qquad\Longleftrightarrow\qquad
    \left[
      \begin{array}{@{}cc|c@{}}
        \Cloze{0.9cm}{a_{11}}&\Cloze{0.9cm}{a_{12}}&\Cloze{0.9cm}{b_1}\\[0.3em]
        \Cloze{0.9cm}{a_{21}}&\Cloze{0.9cm}{a_{22}}&\Cloze{0.9cm}{b_2}
      \end{array}
    \right]
  \]

  \begin{ClozeArea}[34]{24}
    \textbf{Vorwärtselimination:}
    \[
      \begin{array}{r@{\;}r@{\qquad}l}
        &\left[
          \begin{array}{@{}cc|c@{}}
            a_{11}&a_{12}&b_1\\
            a_{21}&a_{22}&b_2
          \end{array}
        \right]
        &\raisebox{-0.65em}{\(\TransformOp{a_{11}Z_2-a_{21}Z_1}\)}\\[1.2em]
        \Longleftrightarrow
        &\left[
          \begin{array}{@{}cc|c@{}}
            a_{11}&a_{12}&b_1\\
            0&a_{11}a_{22}-a_{12}a_{21}&a_{11}b_2-a_{21}b_1
          \end{array}
        \right]
        &\text{Zeilenstufenform}
      \end{array}
    \]
  \end{ClozeArea}
\end{ExampleBox}

\clearpage
\begin{tcolorbox}[
  scriptbox,
  breakable=false,
  borderline west={1.2pt}{0pt}{black!40},
  title={Beispiel~\ref{ex:linear_systems_determinant_from_gaussian_elimination}: Allgemeines $2\times2$-LGS -- Fortsetzung}
]
  \begin{ClozeArea}[34]{22}
    \textbf{Rückwärtseinsetzen:} Für
    $a_{11}a_{22}-a_{12}a_{21}\neq0$ folgt
    \begin{itemize}[leftmargin=2.1em,itemsep=0.7em]
      \item In Zeile 2 einsetzen:
        \[
          (a_{11}a_{22}-a_{12}a_{21})x_2
          =a_{11}b_2-a_{21}b_1
          \quad\Longleftrightarrow\quad
          x_2=\frac{a_{11}b_2-a_{21}b_1}
          {a_{11}a_{22}-a_{12}a_{21}}.
        \]
      \item in Zeile 1 einsetzen:
        \[
          \begin{aligned}
            a_{11}x_1+a_{12}x_2=b_1
            \quad\Longleftrightarrow\quad&
            a_{11} x_1 + a_{12} \frac{a_{11}b_2-a_{21}b_1}{a_{11}a_{22}-a_{12}a_{21}} = b_1 ,\\[0.6em]
            \quad\Longleftrightarrow\quad&
            x_1=\frac{a_{22}b_1-a_{12}b_2}{a_{11}a_{22}-a_{12}a_{21}}.
          \end{aligned}
        \]
    \end{itemize}
    Damit ergibt sich die folgende Lösungsmenge:
    \[
      \mathcal L=
      \left\{
        \left(
          \frac{a_{22}b_1-a_{12}b_2}{a_{11}a_{22}-a_{12}a_{21}},
          \frac{a_{11}b_2-a_{21}b_1}{a_{11}a_{22}-a_{12}a_{21}}
        \right)
      \right\}.
    \]
  \end{ClozeArea}
\end{tcolorbox}

Die Ergebnisse des LGS lassen sich mit der folgenden Definition einer Determinante
bedeutend leichter und übersichtlicher schreiben.

\begin{DefinitionBox}{Determinante einer $2\times2$-Matrix}{linear_systems_two_by_two_determinant}
  Für eine $2\times2$ Koeffizientenmatrix
  \[
    A=\begin{pmatrix}
      \textcolor{THMBlue}{a}&\textcolor{olive!80!black}{b}\\
      \textcolor{olive!80!black}{c}&\textcolor{THMBlue}{d}
    \end{pmatrix}
  \]
  berechnet sich die \emph{Determinante} \EnglishTerm{determinant} von $A$ durch
  \[
    \det(A)
    =\det\!\begin{pmatrix}
      \textcolor{THMBlue}{a}&\textcolor{olive!80!black}{b}\\
      \textcolor{olive!80!black}{c}&\textcolor{THMBlue}{d}
    \end{pmatrix}
    =\begin{vmatrix}
      \textcolor{THMBlue}{a}&\textcolor{olive!80!black}{b}\\
      \textcolor{olive!80!black}{c}&\textcolor{THMBlue}{d}
    \end{vmatrix}
    =\textcolor{THMBlue}{ad}-\textcolor{olive!80!black}{bc}
  \]
  \textbf{Merkhilfe:}
  \enquote{\textcolor{THMBlue}{Hauptdiagonale} minus
  \textcolor{olive!80!black}{Gegendiagonale}.}
\end{DefinitionBox}

\begin{ExampleBox}{Determinante einer $2\times2$-Matrix}{linear_systems_two_by_two_determinant_calculation}
  \begin{enumerate}[label=\alph*),leftmargin=1.8em,itemsep=1.0em]
    \item
      \(
        \det\!\begin{pmatrix}3&2\\1&4\end{pmatrix}
        =\Cloze{6.5cm}{3\cdot4-2\cdot1=10}.
      \)
    \item
      \(
        \det\!\begin{pmatrix}2&-1\\5&3\end{pmatrix}
        =\Cloze{6.5cm}{2\cdot3-(-1)\cdot5=11}.
      \)
  \end{enumerate}
\end{ExampleBox}

\begin{samepage}
Mit der Definition der $2\times2$ Determinante lassen sich die Lösungen des allgemeinen
$2\times2$ LGS aus Beispiel~\ref{ex:linear_systems_determinant_from_gaussian_elimination}
nun sehr übersichtlich schreiben.
\[
  x_1=
  \frac{
    \det\!\begin{pmatrix}
      \textcolor{THMBlue}{b_1}&a_{12}\\
      \textcolor{THMBlue}{b_2}&a_{22}
    \end{pmatrix}
  }{
    \det\!\begin{pmatrix}a_{11}&a_{12}\\a_{21}&a_{22}\end{pmatrix}
  },
  \qquad
  x_2=
  \frac{
    \det\!\begin{pmatrix}
      a_{11}&\textcolor{THMBlue}{b_1}\\
      a_{21}&\textcolor{THMBlue}{b_2}
    \end{pmatrix}
  }{
    \det\!\begin{pmatrix}a_{11}&a_{12}\\a_{21}&a_{22}\end{pmatrix}
  }.
\]
\end{samepage}
